\scanpage{95}
Les arêtes (\deleted{\unsure{supé|supér}} horizontales supérieures) de la forme $\{s,\sigma_\alpha(s)\}$ sont appelées supérieures internes (ou supérieures triviales) -- ce sont celles qui lient des pts d'un $=$ $\widetilde E_\alpha$ $(\alpha\in t_0)$. Exprimons l'axiome 1 pour ces nouvelles flèches, on trouve

\noindent\vrule\hspace{.5em}\begin{minipage}{.89\linewidth}\uline{Prop 2} \quad Soit $s,s'\in\widetilde E_\alpha$ $(\alpha\in t_0)$, $s$ \deleted{\unsure{et}} $s'$ liés (i.e. $s'=\sigma_\alpha s$). Soit $\beta\in t_0$, $\beta\ne\alpha$, et soit $t\in\widetilde E_\beta$. Alors $t$ ne peut être lié simultanément à $s$ et $s'$.\end{minipage}

On appellera \deleted{flèches}\inserted{arêtes} (hor.) sup transverses les \deleted{flèches}\inserted{arêtes} (hor) sup qui ne sont pas internes, i.e. qui \deleted{ne sont} relient un pt d'un $\overline E_\alpha$ à un pt d'un $\widetilde E_\beta$ $(\alpha\deleted{\unsure{\ne}}\unsure{,}\beta\in t_0,\ \alpha\ne\beta)$. Soit $\{s,t\}$ $(s\in\overline E_\alpha,\ t\in\widetilde E_\beta)$ une telle arête, \deleted{\unsure{L'existence de}} \inserted{considérons un} \deleted{le} triangle \deleted{correspondant (par l'axiome 1),} \inserted{$\{s,t,u\}$ la contenant ($\exists$ par} \deleted{\unsure{supposé pour le}} l'axiome 1), mais peu importe pour le moment). Évidemt $u\notin t_0$ (car si on avait $u\in t_0$ on \deleted{devrait} aurait $u=\alpha$ et $u=\beta$, or $\alpha\ne\beta$ !) donc $u\in S\setminus t_0=\widetilde E_\alpha\amalg\widetilde E_\beta\amalg\overline E_{\unsure{\gamma|\alpha}}$. On ne peut avoir $u\in\overline E_{\unsure{a|\alpha}}$, car on aurait alors $u=\sigma_\alpha(s)=s'$ par Prop 1, ce qui contredit prop. 2. Pour la $=$ raison on \deleted{\unsure{doit}} a $u\notin\overline E_\beta$, donc $u\in\widetilde E_{\unsure{\gamma|\sigma}}$, i.e.
% archive sheet number: 94
